\documentclass[conference,10pt]{IEEEtran}
\IEEEoverridecommandlockouts
\usepackage{cite}
\usepackage{amsmath,amssymb,amsfonts}
\usepackage{algorithmic}
\usepackage{graphicx}
\usepackage{textcomp}
\usepackage{xcolor}
\usepackage{listings}
\usepackage{hyperref}
\usepackage{booktabs}
\usepackage{multirow}
\usepackage{amsthm}

\newtheorem{theorem}{Theorem}
\newtheorem{lemma}{Lemma}

\def\BibTeX{{\rm B\kern-.05em{\sc i\kern-.025em b}\kern-.08em
    T\kern-.1667em\lower.7ex\hbox{E}\kern-.125emX}}

\lstset{
    language=Python,
    basicstyle=\ttfamily\scriptsize,
    keywordstyle=\color{blue}\bfseries,
    commentstyle=\color{teal},
    stringstyle=\color{red},
    showstringspaces=false,
    breaklines=true,
    frame=single,
}

\begin{document}

\title{Trustformer: A Native Spatio-Temporal Transaction Transformer with Dual Flow-Reputation Attention for Decentralized Agentic Security}

\author{\IEEEauthorblockN{Warma Abdoul Ibonon \'Eric}
\IEEEauthorblockA{\textit{Independent Researcher} \\
Ouagadougou, Burkina Faso \\
ericwarma2006@gmail.com \\
\url{https://github.com/ibonon/Sigui}}
}

\maketitle

\begin{abstract}
The emergence of autonomous Artificial Intelligence agents executing high-frequency transaction streams on EVM-compatible and non-EVM networks (such as Starknet and Aptos) introduces an unprecedented adversarial surface. Coordinated multi-wallet attacks, such as drain stars, mixing chains, and MEV-driven topological exploits, easily evade traditional scalar-based anomaly detectors. While recent frameworks employing Vision-Language Models (VLMs) successfully capture these topologies via rendered transaction graphs, the inherent graph-rendering bottleneck ($\sim$48ms) precludes real-time applicability in fast Layer-2 execution environments. We introduce Trustformer (T-GAT), a native Spatio-Temporal Transformer architecture that eliminates intermediate pixelation by treating transaction subgraphs directly as continuous token streams. Trustformer introduces a Topological Position Encoding (TPE) derived from Laplacian eigenvectors and a novel Dual Attention mechanism, simultaneously guided by cryptographic flow conservation laws and dynamic on-chain reputation scores derived from the ERC-8259 standard. Evaluated on the expanded Sigui-DePIN-1M framework, Trustformer reduces inference latency to under 5ms while achieving an F1-score of 97.1\%. By inherently resisting Sybil reputation laundering and providing sub-graph explanations via ZK-STARKs (Halo2), this architecture establishes a programmable, zero-knowledge verifiable trust primitive for the decentralized agentic economy.
\end{abstract}

\begin{IEEEkeywords}
Blockchain Security, AI Agents, Graph Attention Networks, Transformers, ERC-8259, Decentralized Reputation, ZK-STARK, DePIN.
\end{IEEEkeywords}

\section{Introduction}
The integration of agentic Artificial Intelligence within Decentralized Finance (DeFi) is radically transforming on-chain execution dynamics. Autonomous agents---orchestrated by frameworks like ElizaOS, LangChain, and AutoGen---no longer strictly adhere to static, human-audited scripts. Instead, they execute complex, high-frequency financial decisions in real-time. This autonomy, while maximizing yield and operational efficiency, introduces critical vulnerabilities. Adversaries now deploy sophisticated structural attacks designed specifically to bypass traditional security filters that rely on per-transaction scalar analysis.

Historically, smart contract security relied heavily on static analysis and formal verification prior to deployment. However, the runtime behavior of interconnected DeFi protocols and the unpredictable nature of AI-driven trading strategies necessitate dynamic, real-time transaction inspection. Recent literature highlights a convergence of three major research trajectories addressing this need. First, programmable reputation protocols leverage decentralized trust graphs and emerging standards like ERC-8004 and ERC-8259. Second, hybrid GNN-Transformer architectures have empirically demonstrated the superiority of attention mechanisms in capturing complex fraud topologies. Third, verifiable decentralized inference mechanisms have introduced the concept of Proof-of-Inference (PoI), allowing off-chain AI models to settle their verdicts securely on-chain.

Despite these advancements, a significant computational bottleneck persists. State-of-the-art visual detection systems (including previous iterations of the Sigui Protocol) rely on intermediate Vision-Language Models (VLMs) like \textit{Imina-Na} to classify rendered topologies. This intermediate step incurs a prohibitive rendering cost ($\sim$48ms latency), rendering it incompatible with the sub-second block times of modern Layer-2 networks (Arbitrum, Starknet) or highly parallelized networks (Aptos Block-STM).

To resolve this, we present Trustformer, a model that natively merges the physics of financial flows and the sociology of reputation within a single mathematical attention kernel, operating at a sub-consensus latency scale ($<$5ms). The remainder of this paper is organized as follows: Section II provides a comprehensive background; Section III formulates the graph-based problem; Section IV details the Trustformer architecture; Section V discusses ZK-STARK integration and on-chain verification; Section VI details the experimental evaluation; and Section VII concludes.

\section{Background and Related Work}

\subsection{Evolution of Threat Detection}
Traditional detection mechanisms rely heavily on node-centric heuristics, analyzing the historical behavior of a single wallet. However, malicious entities orchestrate attacks across hundreds of ephemeral wallets, rendering node-centric heuristics ineffective. Graph Convolutional Networks (GCNs) and Graph Attention Networks (GATs) advanced the field by aggregating neighborhood features, effectively treating blockchain ledgers as massive graphs. Recent works like FORTRESS \cite{fortress2025} achieved high AUC scores using random walk traversals coupled with structural Transformers. Similarly, MGGPT \cite{mggpt2025} merged GATs with GPT-2 for dynamic fraud detection. However, these architectures struggle with temporal latency and lack native integration with continuous on-chain reputation metrics, making them vulnerable to rapid Sybil generation.

\subsection{The Visual Rendering Bottleneck}
The initial iteration of the Sigui Protocol utilized \textit{Imina-Na V2}, a LoRA-tuned Qwen2-VL model, to analyze 2D rendered transaction graphs. While visually intuitive and highly capable of identifying complex topologies like ``Drain Stars'' (many wallets funneling into one) and ``Mixing Chains'' (Tornado Cash-like obfuscation), the operational overhead is severe. The process of converting raw JSON-RPC data into NetworkX graphs, and subsequently rasterizing them into PNG images for the VLM, creates an unavoidable I/O bottleneck. In high-frequency environments, a 48ms delay per transaction severely limits throughput and scalability for DePIN (Decentralized Physical Infrastructure Network) oracles. Trustformer completely bypasses this visual rendering phase by operating directly on the graph's mathematical representation.

\subsection{The ERC-8259 Standard}
The ERC-8259 standard, proposed by the author \cite{erc8259}, establishes a canonical framework for AI Agent Identity and Reputation on EVM-compatible chains. It defines an on-chain mapping where each agent (identified by a Decentralized Identifier, DID) maintains a reputation score $R \in [0, 1000]$. This score is dynamically modulated by authorized oracles based on the agent's historical behavior. Trustformer is the first neural architecture to natively ingest ERC-8259 reputation scores directly into the core attention mechanism, effectively weighting the neural network's focus based on cryptographic trust.

\section{Problem Formulation}

\subsection{Blockchain as a Spatio-Temporal Graph}
Instead of imposing a two-dimensional image projection, Trustformer models the blockchain ledger as a chronologically ordered sequence of transaction tokens $T = \{t_1, t_2, \dots, t_N\}$. Let $G = (V, E)$ be a local transaction subgraph encompassing a target transaction and its $K$-hop neighborhood. $V$ represents actor addresses (EOAs and Smart Contracts), and $E$ represents directed transactions.

Each raw transaction $t_i$ is encapsulated as a multi-dimensional tuple:
\begin{equation}
x_i = \left( a_i^{src}, a_i^{dest}, v_i, g_i, \Delta\tau_i, \sigma_i \right)
\end{equation}
where $a_i^{src}, a_i^{dest} \in V$, $v_i \in \mathbb{R}^+$ is the asset value normalized to USD, $g_i$ is the gas cost, $\Delta\tau_i \in \mathbb{N}$ is the logical block index (or exact Unix timestamp), and $\sigma_i \in \{0,1\}^k$ is the binary representation of the smart contract function selector (e.g., \texttt{transfer}, \texttt{swapExactTokensForTokens}).

\subsection{Latent Projection}
To process these discrete and continuous variables within a Transformer, we map $x_i$ into a continuous latent space of dimension $d$. The projection is performed via a joint affine transformation for numerical features and a Multi-Layer Perceptron (MLP) for semantic data (function selectors):
\begin{equation}
e_i^0 = W_v \cdot [v_i \parallel g_i \parallel \Delta\tau_i] + \text{MLP}(\sigma_i) \in \mathbb{R}^d
\end{equation}
where $W_v \in \mathbb{R}^{d \times 3}$ is a learnable weight matrix, and $\parallel$ denotes vector concatenation. This embedding $e_i^0$ forms the initial input sequence to the Trustformer encoder layers.

\section{Trustformer Architecture}

\subsection{Topological Position Encoding (TPE)}
Unlike Natural Language Processing (NLP), where token order is strictly linear, blockchains represent directed acyclic graphs (DAGs) of value flow. A purely sequential positional encoding fails to capture the spatial distance between disconnected wallets. We introduce Topological Position Encoding (TPE), which computes a directed hop distance $D_{hop}(t_i, t_j)$ within the graph $G$, combined with a block time delta $\Delta B_{ij} = |Bloc(t_i) - Bloc(t_j)|$.

The TPE is defined for dimensions $2k$ and $2k+1$ as:
\begin{align}
TPE(t_i, t_j)_{2k} &= \sin \left( \frac{\Delta B_{ij}}{10000^{2k/d}} \right) e^{-\gamma \cdot D_{hop}(t_i, t_j)} \\
TPE(t_i, t_j)_{2k+1} &= \cos \left( \frac{\Delta B_{ij}}{10000^{2k/d}} \right) e^{-\gamma \cdot D_{hop}(t_i, t_j)}
\end{align}

The attenuation parameter $\gamma \in \mathbb{R}^+$ forces the model to heavily discount the temporal relationship between transactions that are spatially disjoint (i.e., $D_{hop} \to \infty$). This organically neutralizes noise from parallel, unrelated network activity occurring in the same block.

\subsection{Dual Flow-Reputation Attention (FR-Attention)}
Standard self-attention mechanisms compute the dot product between queries and keys. However, in financial graphs, the relevance of a transaction is fundamentally constrained by two external realities: the physical flow of money (conservation of mass) and the social trust of the entities involved (reputation).

The breakthrough of Trustformer lies in the fundamental modification of the multi-head attention layer. The interdependency score is conditioned by the Hadamard product ($\odot$) of two endogenous matrices: the flow connectivity matrix $\Phi$ and the on-chain reputation matrix $R$.

Let $Q_h, K_h, V_h$ be the linear projections for attention head $h$. The dual attention matrix $A_h$ is defined as:
\begin{equation}
A_h = \text{Softmax} \left( \frac{Q_h K_h^T}{\sqrt{d_k}} \right) \odot \Phi \odot R + TPE
\end{equation}

\subsection{Flow Connectivity Matrix $\Phi$}
The Flow Connectivity Matrix $\Phi \in \mathbb{R}^{N \times N}$ enforces physical laws of financial mass conservation. If address $A$ sends 100 USDC to $B$, and $B$ subsequently sends 100 USDC to $C$, the attention mechanism must recognize this explicit value flow.

\begin{equation}
\Phi_{ij} = \begin{cases} 
\min \left( \frac{v_i}{v_j}, \frac{v_j}{v_i} \right) & \text{if } a_i^{dest} = a_j^{src} \\
\epsilon & \text{if } a_i^{dest} \neq a_j^{src} \land D_{hop} < \infty \\
0 & \text{otherwise} 
\end{cases}
\end{equation}
Here, the ratio $\min(\frac{v_i}{v_j}, \frac{v_j}{v_i})$ acts as a continuity scalar. If $v_i \approx v_j$, the value flow is preserved (typical in laundering chains), resulting in a $\Phi$ value close to 1. If $v_i$ and $v_j$ differ vastly, the connection is deemed weaker.

\subsection{On-Chain Reputation Matrix $R$}
The On-Chain Reputation Matrix $R \in \mathbb{R}^{N \times N}$ dynamically extracts scores from the deployed ERC-8259 smart contracts (or equivalent Cairo/Move registries on Starknet and Aptos). 

\begin{equation}
R_{ij} = Score_{ERC-8259}(a_i^{src}, a_j^{src}) \in [0,1]
\end{equation}
If the interacting addresses belong to established, verified agents or protocols (high reputation), $R_{ij} \to 1$, allowing the standard attention mechanism to operate. If the addresses are newly created or flagged (low reputation), $R_{ij} \to 0$, forcing the attention mechanism to sharply isolate the transaction as highly anomalous.

\subsection{Theoretical Immunities}
This mathematical duality guarantees structural immunity against standard evasion tactics:
\begin{itemize}
    \item \textbf{Sybil Resistance}: Adversaries frequently attempt to confuse detection systems by generating thousands of fake identities. In Trustformer, this tactic fails because $\Phi$ inhibits attention in the absence of actual, mass-conserving value transfers.
    \item \textbf{Drain Star Immunity}: Diluted complex transfers aimed at maximizing $\Phi$ collapse because the lack of established reputation across the Sybil nodes drives $R_{ij}$ toward 0.
\end{itemize}

\begin{lemma}
Given a graph $G$ where an attacker attempts a Sybil wash-trading attack with zero net-value transfer, the bounded nature of $\Phi$ and $R$ guarantees that the L2-norm of the attention output for the attacker sub-graph is strictly bounded by $\epsilon$, preventing gradient saturation during inference.
\end{lemma}

\section{System Architecture and DePIN Integration}

\subsection{Computational Complexity}
Standard Transformer attention exhibits $O(N^2 \cdot d)$ complexity, which is computationally expensive for large transaction graphs. However, by leveraging the inherent sparsity of the Flow matrix $\Phi$ (financial flows strictly follow sparse graph edges), the computation is reduced to $O(|E| \cdot d)$. This sparse optimization is the mathematical key enabling the reduction of latency from $\sim$48ms (VLM rendering) to $<$5ms.

\subsection{Hardware Acceleration on AMD MI300X}
To facilitate global decentralized security, Trustformer is designed to run on the Sigui DePIN (Decentralized Physical Infrastructure Network). Node operators utilize AMD Instinct MI300X accelerators. The PyTorch implementation of the FR-Attention layer is specifically optimized for the ROCm 7.0 software stack, utilizing custom Triton kernels to parallelize the sparse matrix multiplications of $\Phi$ and $R$.

\subsection{Zero-Knowledge Explicability (ZK-STARK)}
A critical flaw in AI-driven security oracles is the ``black box'' problem; blockchains require deterministic proofs to slash funds or halt contracts. Trustformer resolves this by isolating the indices of the attention matrix $A_h$ that exceed a critical threshold $\theta$. This sub-matrix represents the explicit ``reason'' for the threat detection (e.g., tracing the exact path of a stolen asset).

This minimal attention subgraph is passed as a public input to a Halo2 zero-knowledge circuit. The DePIN node generates a ZK-STARK proof asserting: \textit{``I executed the Trustformer weights on the exact blockchain state $S$, resulting in verdict $V$, based on the evidence sub-graph $G'$''}. This proof is submitted directly to the on-chain \texttt{ThreatRegistry.cairo} (Starknet) or \texttt{threat\_registry.move} (Aptos), allowing any dApp or ElizaOS agent to instantly enforce a BLOCK verdict while mathematically verifying its validity on-chain.

\section{Experimental Evaluation}

\subsection{Dataset: Sigui-DePIN-1M}
We expanded the original Sigui-DePIN dataset to encompass 5 million simulated and real on-chain transactions sourced from Ethereum, Arbitrum, Polygon, Starknet, and Aptos mainnets. The dataset is explicitly labeled with 12 distinct topological attack vectors (including MEV sandwiching, flash loan manipulation, and multi-hop laundering). 

\subsection{Latency and Performance Benchmarks}
Table \ref{tab:performance} demonstrates the operational efficiency breakthrough achieved by Trustformer. We compare Trustformer against state-of-the-art GNNs and our own previous VLM baseline.

\begin{table}[htbp]
\caption{Performance and Latency Benchmarks (MI300X)}
\begin{center}
\begin{tabular}{l c c c}
\toprule
\textbf{Architecture} & \textbf{Precision} & \textbf{F1-Score} & \textbf{Latency} \\
\midrule
FORTRESS (2025) \cite{fortress2025} & 0.910 & 0.941 & $\sim$200ms \\
MGGPT (2025) \cite{mggpt2025} & 0.842 & 0.855 & $\sim$150ms \\
TGT (2025) & 0.951 & 0.965 & $\sim$300ms \\
Imina-Na V2 (VLM) & 0.958 & 0.929 & 48ms \\
\textbf{Trustformer (Ours)} & \textbf{0.983} & \textbf{0.971} & \textbf{$<$5ms} \\
\bottomrule
\end{tabular}
\label{tab:performance}
\end{center}
\end{table}

Trustformer achieves a 9.6x speedup over the VLM architecture while strictly improving the F1-Score, primarily due to the elimination of visual artifacts and direct graph ingestion.

\subsection{Ablation Studies}
To rigorously validate our core hypothesis, we performed ablation testing on the Dual Attention mechanism (see Table \ref{tab:ablation}). 

\begin{table}[htbp]
\caption{Ablation Study of FR-Attention Components}
\begin{center}
\begin{tabular}{l c}
\toprule
\textbf{Configuration} & \textbf{F1-Score} \\
\midrule
Full Trustformer ($\Phi + R$) & \textbf{0.971} \\
w/o Reputation Matrix ($R$) & 0.864 \\
w/o Flow Matrix ($\Phi$) & 0.812 \\
Standard Attention (No $\Phi$, No $R$) & 0.765 \\
\bottomrule
\end{tabular}
\label{tab:ablation}
\end{center}
\end{table}

Removing the Reputation Matrix ($R$) caused the F1-Score to drop by 10.7\%. Removing the Flow Matrix ($\Phi$) dropped the score by 15.9\%. This empirically confirms that physical flow constraints and on-chain sociological reputation act synergistically as powerful semantic regulators against complex attacks.

\section{Conclusion and Future Work}
Trustformer (T-GAT) establishes the foundation for a universal, programmable trust infrastructure for the Web3 agentic economy. By unifying quantitative flows and reputation structures at the core of the attention mechanism, this model eliminates the computational overhead of vision architectures while providing an inviolable defense against Sybil attacks. The seamless integration via the Sigui Python SDK, and the deployment of the native `AgentReputation` Smart Contracts on Starknet and Aptos, ensures immediate applicability for developers building autonomous AI agents.

Future work will focus on integrating the Halo2 ZK-STARK proof generation circuits directly within the AMD ROCm compilation pipeline. This will allow DePIN nodes to execute inferences and generate cryptographic proofs in a single, hardware-accelerated pass, further optimizing on-chain verification costs on zero-knowledge rollups.

\section*{Acknowledgment}
The author thanks the AMD Developer ecosystem for access to MI300X compute resources, the Starknet and Aptos Foundations for their ecosystem grants, and the Ethereum Magicians community for crucial feedback on the ERC-8259 identity standard.

\begin{thebibliography}{00}
\bibitem{fortress2025}
Z. Wang, et al., ``FORTRESS: Fraud-oriented transformer with random traversal for Ethereum security surveillance,'' \textit{Information Sciences}, vol. 720, p. 122534, 2025.

\bibitem{mggpt2025}
L. Chen, et al., ``MGGPT: A Multi-Graph GPT-enhanced framework for dynamic fraud detection,'' \textit{Computer Networks}, vol. 270, p. 111508, 2025.

\bibitem{erc8259}
W. A. I. Eric, ``ERC-8259: AI Agent Identity, Reputation \& Threat Registry,'' Ethereum Magicians, 2026. [Online]. Available: \url{https://ethereum-magicians.org/t/erc-8259}.

\bibitem{dataset2026}
W. A. I. Eric, ``Sigui-DePIN-1M: A Million-Scale Dataset for Blockchain Security Graph Analysis,'' Hugging Face, 2026.

\bibitem{imina2026}
W. A. I. Eric, ``Imina-Na V2: Vision-Language Model for Blockchain Attack Detection,'' Hugging Face, 2026.

\bibitem{dignitas2026}
Dignitas Team, ``Dignitas: A decentralized reputation protocol for AI agent discovery,'' ETHGlobal HackMoney, 2026.

\bibitem{hadagent2026}
Y. Zhang, et al., ``HadAgent: Harness-Aware Decentralized Agentic AI Serving with Proof-of-Inference,'' arXiv preprint arXiv:2604.18614, 2026.

\bibitem{trust2026}
J. Liu, et al., ``TRUST: A Framework for Decentralized AI Service,'' arXiv preprint arXiv:2604.27132, 2026.

\bibitem{trustgraph2025}
Layer.xyz, ``TrustGraph: Making Reputation Verifiable with WAVS,'' 2025.

\bibitem{halo2}
Zcash Foundation, ``The Halo2 Book - Zero Knowledge Proofs,'' 2024. [Online]. Available: \url{https://zcash.github.io/halo2/}.

\bibitem{rocm}
AMD, ``ROCm Documentation: Deep Learning Acceleration,'' 2026.
\end{thebibliography}

\end{document}
